\section{Discussion}\label{sec:4}

\subsection{Two regimes, one object each, and a gap between them}\label{sec:4.1}

The corpus is not silent about police robots. It contains 1,014 authorising provisions, of which the predicate returns 91 and the audit confirms 41, and the instruments that produce them

are recent, specific and enforceable. What the measurement shows is that those 41 bindings attach to three objects, and that the object a deployment decision turns on is none of them.

The first object is the machine as an article of commerce. Airworthiness, type certification, conformity assessment and market access account for 6 of the 41 audited bindings, assessed before sale, by a body external to the buyer, against thresholds fixed in advance. That is a minority, and an earlier version of this paper had it as the majority; the classification of Section 3.3 measures 35 of the 41 as conditioning an operating permission instead. The correction matters for how the gap should be described. Law does not confine itself to the point of sale and leave use alone. It conditions use heavily, on pilot licences, operator certificates, aerodrome certificates, operational authorisations and the airframe's own compliance, and the property it conditions use on is never a sensing performance. The second object is the purpose of the police act. The Law Enforcement Directive, the AI Act's law-enforcement provisions and the municipal and state equipment-governance instruments are precise about why a capability may be used, who must approve it and what follows from a breach. Between them lies the proposition that the reading on which an officer acts has been shown to be adequate for the conditions of use. No audited binding occupies that space. The three instruments that regulate police conduct return none in their 30 authorising provisions, and the two horizontal product instruments return three in 139, all rejected by the audit.

The third object is the one this paper discovered by failing to refute itself, and it is the smallest of the three: a single provision, 14 CFR \S{}89.110(b), conditions permission to operate on an accepted declaration that the aircraft reports its own geometric position to within 100 feet with 95 per cent probability (Section 3.4). Most of the audited bindings condition operating rather than selling, so that is not what distinguishes it; what does is that it is the only one of the 23 aviation bindings that survives the police carve-out. Its object is the machine's account of its own location. This is a performance property, it is stated numerically, it is tied to a permission determinately, and it concerns the aircraft as a thing to be found in the sky rather than as an instrument that produces a reading an officer will act on. The gap this paper describes is therefore not a gap between regulation and silence. It is a gap between the properties of a police robot that law has learned to condition permissions on and the property a police decision actually rests on, and those two sets intersect in exactly one provision concerning neither the robot's sensing nor the officer's use of it. The police-conduct zero rests on 30 provisions and is reported in Section 3.3 as a reading of their text rather than a statistical contrast.

The regulatory-design literature explains why the division falls where it does. Kaplow's distinction between rules and standards turns on when the content of the law is fixed \citep{Kaplow1992Rules}, and Diver's account of optimal precision adds that precision is costly and is therefore bought where the rulemaker can bear the cost of being specific \citep{Diver1983Precision}. Specifying in advance what sensing performance suffices for a police determination requires the rulemaker to know the mission, the analyte, the environment and the consequence. Specifying that a permission depends on a purpose and a prior authorisation requires none of that. Coglianese and Lazer's account of management-based regulation describes what is written instead: the regulated party must plan, assess and document, and the regulator supervises the process \citep{CoglianeseLazer2003}. Such supervision is a reasonable response to heterogeneous technology, and it is precisely the form our predicate records as indeterminate.

\subsection{The asymmetry runs between certification and conduct, not between aviation and the ground}\label{sec:4.2}

The pooled comparison between aviation at 10.8\% and the ground corpus at 2.5\% invites the obvious explanation, that aviation certification is a century old and ground robotics is not. We

cannot offer that explanation, because we cannot defend the comparison it would explain. The ground category pools two bodies of law that have no reason to behave alike, and 95\% of its denominator is human driver licensing (Section 3.3). Splitting it does not rescue a different comparison either: the Korean robot instruments supply 30 authorising provisions and 1 audited binding, and a cell that size cannot distinguish the aviation rate from the road-traffic rate in either direction.

The densest contrast is aviation certification against road-traffic law, 26 of 241 against 14 of 574, and it is large in the descriptive sense established in Section 3.3 rather than in a tested one. That is a contrast between a regime built to condition a permission on a demonstrated property of a machine and a regime built to condition a permission on a demonstrated property of a person. The ground bindings are driving examinations and aptitude tests; the aviation bindings are airworthiness and conformity assessment. Whether robot certification behaves like the first or the second is a question this corpus is too small to answer, and we leave it open rather than answering it from two events.

Two explanations of the aviation result itself have to be set aside. The first is a language artefact: an earlier version of this paper reported a Japanese zero, and one might read the asymmetry as an instrument that works in English aviation prose and fails elsewhere. The zero was a lexicon failure, and it has since been repaired (Section 3.3); the Act now yields 8 audited bindings and the Korean Road Traffic Act yields 5, both in languages the instrument was accused of failing in. The asymmetry survives the repair at a smaller effect size. The second is sample size in the other direction: the police-conduct instruments return nothing from 30 provisions, which is not separable from the aviation rate at that denominator, so what we say about them rests on reading their text and not on a test.

The surviving contrast is narrower than the version of this section we first wrote, and we record the difference because it matters to how the paper should be used. We claimed a contrast between regulatory domains. What the data support is a contrast between two specific bodies of law, and a demonstration that our own pre-registered domain partition pooled instruments it should not have.

Directive 2014/32/EU conditions the marketing of a measuring instrument on a stated maximum permissible error, and lists the instruments it governs: water, gas, electricity and heat meters, fuel dispensers, automatic weighing instruments, taximeters, material measures, dimensional instruments and exhaust gas analysers. The legal form that would bind a police sensor to a demonstrated accuracy has existed in its present shape since 2014 and in its predecessor since 2004. It was not withheld because the form had not been invented. It was withheld because the scope was drawn by enumerating instrument categories, and law-enforcement measurement was not one of them.

The same point holds inside aviation, with a caveat the audit forces on us. Part 89 was made in 2021 and requires a standard remote identification unmanned aircraft to broadcast a specified message set against an FAA-accepted declaration of compliance. A technical requirement on an unmanned aircraft was therefore available to be written five years ago, and one was written; what it requires is that the aircraft announce who it is. Our audit does not count it as a determinate binding of an authorisation, for the reasons given in Section 3.4, and the argument here does not need it to be one. The point is about what the drafter chose to specify when specifying was plainly possible.

\subsection{A carve-out is a change of object, not a withdrawal of law}\label{sec:4.3}

The aerial comparison isolates the asymmetry inside a single mission. Twenty-three audited bindings apply to a civil unmanned aircraft across the two jurisdictions; one applies on the police route. The EU half of that result is close to analytic and we say so in Section 3.4; the half that measures something is the US one, where Part 89 remains in force for a police operator and one of its 17 authorising provisions binds. Section 3.4 sets out why the retained set is Part 89 alone and why our deletion, which also retained parts of Part 107 and the advisory circular, erred in the direction that makes the police route look better covered. The usual description of a publicaircraft carve-out is that police operations are less regulated. The measurement says something more specific: the police operator is subject to the same purpose rules, the same data rules and the same equipment-governance rules, and is released from almost the whole of the layer that would have required a demonstration that the machine meets a stated physical threshold before it flies over a person. The release is not total, and the exception is informative in a way a clean zero would not have been.

The carve-out has a coherent rationale. Public aircraft operations are self-certified so that an agency can act in an emergency without waiting for a certification process built for commerce. The finding is about what that rationale costs and where the cost lands. It lands on the one layer the purpose rules cannot substitute for, because a determinate rule about why a robot may be used says nothing about whether its readings are good enough to act on.

The survivor sharpens the point rather than blunting it. It requires that the aircraft's reported geometric position be accurate to within 100 feet of the true position with 95 per cent probability, and conditions permission to operate on an accepted declaration to that effect (Section 3.4). Two features of it matter. The first is that it demonstrates the form is available on an operating permission and not only on a market one: a rulemaker conditioned flying, not selling, on a numeric accuracy figure, held the operator rather than only the producer to it, and did so for a class of aircraft that police agencies fly. The objection that determinate performance conditions cannot be written for a police robot because the mission is heterogeneous is answered by a provision that already exists. The second is what the threshold governs. Remote identification serves third-party oversight: it lets a person on the ground, or another airspace user, learn what is overhead and who is operating it. A position accurate to 100 feet is what makes that oversight work. It is not a condition on decision quality, and it is not pointed at the reading an officer acts on.

So after the exemption the single technical property of a police drone that law determinately conditions is how accurately it reports where it is, and not how well anything it carries measures the world. An earlier version of this paper wrote that sentence as an inference from a zero. It is now a measurement with one event in it, which is a stronger thing to have: a zero is consistent with a legal system that has not found the form, whereas a single surviving binding shows the form was found, applied to an operating permission, and pointed somewhere else.

\subsection{What a binding contains, and the element the one working example omits}\label{sec:4.4}

The Korean outdoor mobile robot certification is the only place in the ground corpus where a robot's certified condition determinately conditions a legal operating right, and it is instructive for both what it contains and what it omits. It binds a named evaluation to a specific operating right, with a change in the certified configuration re-opening the evaluation \citep{ROKIntelligentRobotsAct}. The evaluation itself is set out in an implementing ministerial notice, distributed as a proprietary word-processor file and therefore outside the machine-readable

corpus; we read it clause by clause and record the reading rather than adjudicating it \citep{MOTIENotice2023211}. Its sixteen articles run from application intake through document examination, product examination, partial or full exemption from examination, disposal of the result and issue of a numbered certificate, with the processing period fixed and an objection procedure provided. The examining body is named and nothing in the transition is at the discretion of the deploying party.

Across those sixteen articles and four appendix tables the notice contains no operator or supervisor proficiency requirement of any kind: the words for operator, qualification, training and licence do not occur in it. What our reading recorded is that absence. It did not adjudicate whether the technical criteria in the appendix tables set a threshold on a sensing or measurement function, and we state the gap rather than close it, because this is the one place in the corpus where our own defeater condition could be met: a numbered certificate that a legal right depends on, whose criteria we have not put through the predicate. The notice is distributed in a proprietary format and is outside the machine-readable corpus, so we cannot adjudicate it under the frozen rules without changing what the corpus is. A reader who obtains it and finds a sensing threshold in those tables has refuted the negative claim for the Korean ground case, and we would rather say so than leave the possibility unstated. The most completely specified binding available in the ground corpus fastens the machine and its control system to the operating right and leaves the human competence that mediates between them outside the certificate. Two of the twelve causal constraints derived in Section 2.2 concern exactly that competence: whether operator qualification is demonstrated against the degraded conditions the mission presents, and whether the controls that must remain available inside a contaminated zone can be actuated in the protective equipment the zone requires.

A second limit is structural. The Korean certificate is indexed to a platform and one operating context, while police missions vary in the dimensions it holds fixed. The same ground robot performs barricade reconnaissance under tactical command and hazardous sampling under evidentiary constraint, and what must be demonstrated differs. What would close the distance is an object indexed to the mission rather than the platform: a record, for a specific deployment, of which evaluation was performed on which sensing function for which analyte and range, what the result was, who accepted it, and what change of condition re-opens it. The components exist separately. A capability certificate answers whether the device met defined thresholds; an operator credential answers whether the user can perform the tasks; only the object that fastens a stated technical result to a specific authorisation is missing, and it is the one that converts a concept of operations from a narrative into an auditable record.

\subsection{Why this cannot be closed by testing alone}\label{sec:4.5}

A testbed can produce the evidence that makes a determinate threshold possible, and that is the mechanism by which an indeterminate requirement becomes a determinate one. It cannot produce a forum competent to read the evidence, and it cannot confer standing to raise the matter. Selbst and colleagues describe the failure mode in which a sociotechnical problem is solved within the boundary of the technical system while the institutional arrangements that would make the solution meaningful are abstracted away \citep{Selbst2019Abstraction}. The pattern measured here is that abstraction written into law: the technical layer and the accountability layer are each determinate about their own contents, and the interface between them is where determinacy stops.

The regulatory timetable offers a way to test this prospectively rather than argue it. The AI Act's Annex III obligations for law-enforcement systems apply from 2 December 2027 \citep{EUAIAct2024,EUDigitalOmnibus2026}, and the EU machinery regulation from 20 January 2027 \citep{EUMachinery2023}. Both convert obligations our predicate currently records as indeterminate or absent into binding requirements on a known date. Re-running this measurement after those dates, with the rules fixed in advance and published, would show whether the gap closes.

\subsection{Limitations}\label{sec:4.6}

One in three of the provisions the predicate returns is a binding. Precision is 0.451 with a Wilson 95\% interval of [0.352, 0.553], and it is not uniform: 0.481 in aviation, 0.324 on the ground and 0.000 across the cross-cutting instruments (Section 2.7). This is the dominant limitation and it governs how every count in this paper should be read. A returned count is an upper bound on binding and nothing more, which is why we report the audited count beside it everywhere and rest the argument on the second. It also sets a floor on what the method can currently deliver without a reader: at this precision, an unaudited count of determinate bindings is not a usable measurement of a single corpus, although the audited comparison between two routes through one corpus remains informative because both routes are audited under the same labels.

The instrument was not symmetric across languages, and the Japanese zero was ours, not the law's. For most of this study the Act returned 11 provisions and the audit rejected all 11, while the Korean Road Traffic Act yielded five bindings under the same rules. The cause was identifiable and we named it before we repaired it: the evaluation lexicon carried the Korean term for a licensing examination and not the Japanese one. We first reported the gap rather than repairing it, on the ground that the ruleset had been frozen. That was the wrong call, and the tier-0 sample of Section 2.10 is what showed it: the omission was a class rather than a term, and leaving it stood only because its direction happened to suit us. The lexicons are now symmetric in both directions, hedges included, and the Act returns 20 provisions of which 8 are bindings. The residual limitation is narrower but real. Symmetry was established against Korean, so a concept absent from both lexicons is still absent from both, and no third-language check exists.

Paragraph granularity and character windows cost bindings in both directions. A determinate condition and a hedge that govern different clauses of the same paragraph are treated as one unit, which removes bindings; relaxing that rule entirely raises the corpus count from 91 to 116 and adds no audited binding. The determinate patterns scan a character window rather than a clause, which is what lets an exemption read as a binding and what defeats the obvious repair (Section 2.7). Annexes are held as single nodes truncated at 6,000 characters, so most annex text is not in the corpus at all rather than merely unresolved to an item (Section 2.4); this cuts both ways, since an annex opening that happens to contain an evaluation term lets any provision referring to it reach one. All three are consequences of working on published text with pattern rules instead of a clause parser, and a clause parser is the single change that would most improve this instrument.

The analyst-dependent step we could not make reproducible. Translating a derived safety constraint into search terms for four legal systems is done by a person, and when we tested it against a second term-writing pass made blind from the same rules, outcome-class agreement was 5 of 12 and Cohen's $\kappa$ was 0.184. We therefore report no constraint-level counts. Everything above is computed over the full population of authorising provisions, where no such translation occurs. The binding audit was checked the same way and is more stable: agreement over all 74 labelled records is 51, $\kappa$ = 0.448, and outside one ambiguous sentence pattern 51 of 56, $\kappa$ = 0.835. The second figure is conditional on removing the records where the disagreement

concentrated and is not a reliability estimate; we give it because it localises the ambiguity, not because it replaces the first. Neither check is agreement between two people. Both repeat a procedure under one investigator, using a language model as the instrument, so they bound stability under repetition and not inter-rater reliability, and we claim only the former.

The lexicon leaks, and the leak is bounded rather than eliminated. The validated falsenegative rate in the population at risk is 12.6\% [4.7, 29.5], projecting to 23 to 145 missed bindings. Carried onto the audited counts with numerator and denominator drawn from the same population, that implies a recall of 0.40 [0.22, 0.64] corpus-wide and, on the police route where the key negative claim lives, 0.26 [0.13, 0.48] (Section 3.4). Those intervals treat provisions as the sampling unit. The sample is clustered by instrument in the same way the corpus is, and correcting for that widens the rate to [2.5, 45.0] per cent, the projection to 12 to 221, and the corpus-wide recall to [0.16, 0.77] (Section 2.6). The wider pair is the honest bound on this claim and the narrower pair is what the stratification was designed to estimate; we give both because the difference between them is larger than anything else in this paragraph. An earlier version of this paper reported the complement of the false-negative rate, 88.6\%, as though it were recall; it is not, and the correction moves the instrument from one that finds most bindings to one that finds a minority of them. None of the observed false negatives concerns sensing or measurement performance, but the labels came from a single pass, and the projection assumes the sample is representative of a population the Japanese gap above shows to be heterogeneous. The exposure is narrowed not by the projection but by the vocabulary sweep of Section 3.6: a missed binding on sensing performance would have to be written without any of seventeen performance concepts, each expressed in all three languages of the corpus. We cannot exclude one, and we do not claim to.

Sixteen provisions carry the whole predicate inside one paragraph and are still scored indeterminate. A provision in which the authorisation, the determinate operator and the evaluation term all occur in the same node is where the predicate should be most reliable. Eighty-eight provisions are of that form, and we counted the ones the predicate still declines after the lexicon symmetrisation: there are 16, and all 16 are tier 1, held there by a hedge in the same paragraph. Twelve of the hedges are the bare modal may or its Korean and Japanese equivalents, used in the sense of is permitted to rather than as a hedge on the condition; the other four are as appropriate twice and two forms of shall take into account, one English and one Japanese. Section 89.115(a) is among the 16, and it is the sister provision of the one binding that survives the carve-out. We report the count rather than resolve it, because resolving it means deciding clause by clause whether a modal governs the permission or the condition, which is the clause-parser problem above and not a lexicon entry; but a reader should know that the same mechanism that hid \S{}89.110(b) for the history of this study is still holding twelve provisions at tier 1, and that the audited count of 34 is therefore a lower bound in a way the lexicon projection does not separately capture.

One part of that mechanism is asymmetric across languages and we had not said so. The English hedge entry for may carries an exception: where a limiting construction follows in the same sentence, the modal is not counted as a hedge. That exception is what recovered \S{}89.110(b), the one binding that survives the carve-out. The Korean and Japanese entries for the same permissive modal carry no corresponding exception. Fifteen Korean and Japanese authorising provisions, two Korean and thirteen Japanese, hold a limiting construction and the bare modal in the same node; one reaches tier 2 already, nine are at tier 1 and five at tier 0, and of the nine at tier 1 exactly two, Articles 90(1) and 97-2(1) of the Japanese Road Traffic Act, already carry an evaluation term and a determinate operator and are held below tier 2 by the modal alone. We do

not extend the exception, for the reason given above: deciding whether a modal governs the permission or the condition is the clause-parser problem, and doing it for Japanese and Korean by pattern would be the kind of rule we would be tuning against the result. But the repair that produced this paper's surviving binding has been applied in one language and not the others, and a reader weighing the Japanese count should know it.

The denominator is contaminated and we measured it rather than assumed it away. An audit of 79 authorising provisions puts the precision of the authorising population at 0.494 [0.386, 0.602], and the contamination is uneven: heaviest among the cross-cutting instruments and lightest on the ground (Section 2.8). Correcting for it raises every rate and widens the contrast the paper rests on, so the nominal rates are the conservative statement and they are the ones quoted throughout. The first draw found three penalty provisions that the exclusion patterns missed, all in the Korean and Japanese material; the second draw found none, so the exclusion is not shown to be complete and we report that rather than repair it against a sample already seen.

A verification suite can confirm consistency and miss construct validity, and ours did. The determinacy lexicon suppressed may . . . only if, the strongest determinate construction in the English material, for the whole of the study's history, and every internal check passed throughout because the checks test whether the reported numbers agree with one another rather than whether the predicate detects what it claims to detect (Section 2.10). The defect was found by taking the paper's own declared defeater seriously and searching for a provision that would refute us. This is a limitation of our verification practice and of the general practice of validating automated legal measurement by internal consistency, and the lesson generalises past this paper: a stated defeater, executed as a search, found in one pass what the consistency suite described in the data statement had not.

Four layers that bind police robots are invisible to a corpus of published provisions. The first is the individual authorisation: in the United States a police unmanned aircraft flies under a Certificate of Waiver or Authorization whose special provisions, not the text of Part 107, carry the operative conditions, and the advisory circular that defines the carve-out says in terms that an agency may elect to operate under the civil rules instead. The second is procurement: a performance threshold written into a tender specification or a testing-laboratory accreditation requirement binds a supplier without appearing in any statute. The third is national legal metrology. Article 3 of Directive 2014/32/EU is headed Optionality, and Article 3(1) provides that Member States may prescribe the use of measuring instruments for measuring tasks where they consider it justified for reasons including public interest, public safety and public order, Article 3(2) requiring only that a Member State which does not prescribe such use communicate its reasons. The decision to impose metrological control on such a task is therefore national and optional, so the enumeration we report in Section 3.5 bounds Union-level requirements rather than the whole set of instruments under metrological control. A fourth, the law of evidence, conditions the use of a measurement rather than the authorisation of the act, which this predicate is defined not to see. None of these is a defect in the measurement. They bound what the measurement is about, and a reader who concludes from this paper that nothing constrains police sensing performance has drawn a conclusion the design cannot support. The narrower claim it does support is that the provisions which authorise, as published, do not carry that constraint.

The corpus is bounded and the claim is bounded with it. Sixteen instruments in four national and two sub-national jurisdictions is a defensible universe for this question and it is not the world. The Korean Measurement Act could not be retrieved in machine-readable form and is the most consequential known omission, since it is the national instrument that would implement metrological verification for Korean police equipment, and it sits in precisely the national layer

the preceding paragraph names.

This study measures law, not practice. A determinately bound provision may be complied with poorly, and an unbound one may be covered in practice by a conscientious standing operating procedure. What the predicate reports is whether an authorisation fails when an evaluation is not passed, which is a statement about enforceability. Establishing the relationship between that and what agencies do requires field observation this design cannot supply.

\section{Conclusions}\label{sec:5}

Whether a police robot may be deployed is a question about an interface: a measurement is taken, a result is produced, and a person with authority decides something with consequences for other people. This study asked where law makes that authorisation depend on a demonstrated property of the machine, and answered with a predicate decided from the text of 16 instruments, applied to all 1,014 authorising provisions they contain, under rules frozen before the answer was seen and then audited provision by provision in two labelling passes whose limits we state in the Methods.

The audit is what the answer rests on. The predicate returns 91 provisions and 41 survive, a precision of 0.451 [0.352, 0.553], so a returned count is an upper bound and every figure below is the audited one. A second audit, of the denominator, puts its precision at 0.494 [0.386, 0.602], so the rates are reported in two coordinates and the nominal pair is the conservative one. Thirty-seven of the 41 bindings are legible in a single paragraph read alone; following cross-references adds four bindings and fifteen false positives, so precision falls from 0.514 to 0.451, which reverses the expectation that motivated building the measurement as a graph in the first place and which we report as a negative result rather than repair. Aviation certification converts 10.8\% of its authorising provisions into audited bindings against 2.4\% for road-traffic law, the contrast in the corpus with the largest event counts; the robot and police-conduct instruments supply 30 provisions each and are reported without a statistical claim.

Six of the 41 bindings attach to the machine as an article rather than to the mission. All six are pre-market certification in the European basic act, assessed before sale by a body external to the buyer against thresholds fixed in advance. The other 35 condition operating it (Section 3.3). Applying the police and public-aircraft carve-outs as a graph operation, holding corpus and rule constant, removes 25 of the 26 audited bindings that a civil unmanned aircraft carries. The European half of that is close to analytic. The American half is the measurement: Part 89 stays in force for a police operator, 17 authorising provisions, and exactly one of them binds.

The one surviving provision is the finding this paper did not expect, and the reason its claim is narrower than the one we set out to test. Section 89.110(b) of 14 CFR permits operating a standard remote identification unmanned aircraft only if its serial number appears on an accepted declaration of compliance, and the declaration attests, four provisions on, that the aircraft reports its own geometric position to within 100 feet of the true position with 95 per cent probability. It is a determinate operator, a named evaluation and a numeric threshold, in force on a police operator. Law has therefore learned to condition a police robot's permission to fly on a demonstrated performance property, and the property it chose is the accuracy with which the machine reports where it is. Nothing in force on that route conditions any authorisation on the accuracy of what the machine senses: not a chemical identification, not a thermal or acoustic reading, not a visual

match. An independent sweep of seventeen performance concepts across all 5,758 nodes finds no such provision in any of the sixteen instruments, on either route.

The finding is a negative one about sensing and its boundaries are part of it. At a precision of 0.451 the instrument cannot deliver an unaudited count worth reading; its recall, carried from the validated false-negative rate onto the audited counts, is near a third corpus-wide and far lower on the police route, so the negative claim is bounded by what the predicate can see rather than by what the law contains. The Japanese zero was a defect in our lexicon rather than a fact about Japanese law, and repairing it cost us a third of the effect size we had reported. The determinacy lexicon also carried, for the whole history of this study, a defect that suppressed the very construction \S{}89.110(b) is written in, and every internal consistency check passed while it did; we found it by executing the paper's own declared defeater as a search, and we report that as a caution about validating automated legal measurement by internal consistency. And the layers that do bind police robots in practice, among them the individual authorisation, procurement and national legal metrology, are invisible to a corpus of published provisions.

The claim that survives all of this is narrow and testable: among the provisions that authorise, as published, the only one that reaches a police robot and conditions its permission on a demonstrated performance conditions that permission on the robot's account of its own position, and none conditions it on the quality of the readings an officer will act upon. This is a claim about what certification regimes were pointed at, not about how young a technology is.

The legal form that would impose the missing condition is neither missing nor immature. Directive 2014/32/EU has conditioned market access on a stated maximum permissible error since 2014 and lists ten categories of instrument to which it applies, none of them a law-enforcement measuring instrument; its own provision fails our audit, so we report this as a reading of the Directive and not as a measured positive control. Part 89 shows the same form working on an operating permission rather than a market one, which is the harder case and the one closer to a deployment decision. Closing the distance therefore does not require inventing an assurance mechanism. It requires pointing one that demonstrably works, on both a market and an operating permission, at the measurements a police determination rests on, and building the accountability forum that could read what it produces.
